
\begin{abstract}

As LLM agents integrate with increasingly large tool ecosystems, tool routing---determining which tools to invoke for a user query---becomes a critical bottleneck.
Existing approaches formulate routing as direct tool selection, requiring the model to simultaneously identify relevant tools, infer multi-step compositions, and ensure structural validity in a single unconstrained decision.
As catalogs grow, this formulation leads to missed compositions, hallucinated invocations, and outputs whose validity cannot be verified without execution.

We argue that this difficulty stems from conflating two distinct reasoning tasks that should be solved separately: \emph{semantic reasoning} (what domain entities does the user refer to?) and \emph{compositional reasoning} (what tool chain connects them?).
We propose \textbf{Typed Composition Routing (TCR)}, a reformulation that decomposes these concerns.
Tools are modeled as typed transformations between entity types, forming a directed composition graph.
The LLM handles only semantic reasoning---predicting source and target entity types---while deterministic graph search handles compositional reasoning, recovering valid tool chains with structural validity guaranteed within the typed composition graph.

Across four language models, five tool domains (54--170 tools), and 140 benchmark queries, this reformulation improves F1 in all 20 model--domain combinations (average $\Delta$F1 = +0.32) without training or fine-tuning.
A production-scale evaluation on the Ansible Automation Platform MCP Server---a real-world system exposing 1,060 API operations across four enterprise services---confirms that the decomposition extends beyond benchmark-sized domains.
The typed composition graph is derived automatically from OpenAPI specifications without manual curation.
At this scale, TCR achieves F1 up to 0.91 while reducing prompt tokens by 94\%.
An analytical recall decomposition separates end-to-end performance into entity type prediction accuracy and graph reachability, attributing routing failures to one component or the other.

\end{abstract}
